% chapter5_conclusions.tex
\chapter{Conclusions \& Practical Implications}
\label{ch:conclusions}

\section{GEO Is Mostly SEO}
\label{sec:geo-is-seo}

Commercial prompts trigger web search over half the time, three
times the baseline rate (Figure~\ref{fig:trigger-rate}), making
retrieval the default for product recommendations. Across 711
runs spanning both models and three deployment tiers, 78--84\% of
LLM citations trace back to a conventional search index
(Section~\ref{sec:global-overlap}). We established
that LLMs do look deeper than human users but citations from beyond
rank~10 drop off sharply (Section~\ref{sec:page-dist}) which is
the backbone of our argument that GEO is mostly SEO.

Within the ranked results, models show consistent structural
preferences. For listicles, pages with tables, numbered lists, and
bullet points are over-selected relative to the SERP baseline. We
saw a similar selection pattern on product pages for numbered lists
(Section~\ref{sec:selection-drift}). But these features were already
standard SEO advice before GEO became a term: they help human
readers scan and compare products, and search engines have long
rewarded content that serves readers well. The fact that LLMs also
prefer them is not necessarily a new optimization surface. We showed that Gemini pre-filters for
recency through year-token injection in fan-out queries, and GPT Plus
shows genuine freshness preference at the citation stage
(Section~\ref{sec:freshness-subsection}). Keeping content current
has been good SEO practice and it is now also a good GEO
practice.

Listicles are the primary fuel for product recommendations and they
are also the page type that predominantly reaches the model through
SERP results. A listicle must rank to be retrieved in the first place. A brand could, in theory, craft a ``top~10 AI tools''
listicle that games our findings, burying real competitors at the
bottom of the list, omitting them entirely, or only listing soft
alternatives. But as long as search engines still weigh engagement
signals like bounce rate and time on page, a listicle that readers
find unhelpful will not rank or will not maintain rankings long term.
An unranked listicle is invisible to the grounding pipeline
regardless of how well it is structured or how well it plays the LLM
manipulation game. GEO optimization relies on the SEO optimization
loop: a listicle that serves human readers earns the rankings, and
the LLM citations follow from that.

Our cross-language fan-out analysis
(Section~\ref{sec:localization-bias}) adds a localisation dimension
to GEO. Due to the English bias in LLM fan-out queries
\parencite{rudzki2026english}, a foreign-language business that
lacks English-language pages will miss a substantial share of
retrieval opportunities, since 43\% of all fan-out queries are
issued in English regardless of the prompt's language.

The workflow that GEO platforms are beginning to offer (identifying
the prompts a brand appears in and tracking visibility over time)
needs to be complemented by an understanding of fan-out queries. Once
a brand maps the fan-out queries that LLMs generate for prompts in
their niche, they can run a keyword gap analysis on those queries,
which mirrors traditional SEO gap analysis.

The one partial exception is the invisible-link channel: 16--22\% of
citations point to URLs that do not appear in either Bing or Google
for the corresponding query (Section~\ref{sec:invisible-links}),
concentrated in Wikipedia, Reddit, app stores, and major tech
publications. These sources enter the response through pathways that
appear to bypass search results. They give brands a form of visibility that is durable and difficult for competitors
to displace,
as long as the current grounding workflow remains in place.

The entire chain (query generation, retrieval, source selection,
product extraction) is grounded in the same web infrastructure
that SEO has optimized for decades. A brand that already runs
rigorous SEO gap analysis is, in practice, already covering most of
the GEO gap.

\section{Outlook}
\label{sec:outlook}

One finding we did not anticipate was how differently the two ChatGPT
tiers behave. There had been reports that OpenAI draws on Google's
search data alongside Bing
\parencite{theinformation2025openaiserpapi}, and our data is
consistent with this: Business citations overlap with Google at
only 27.8\% while Plus reaches 64.6\%.
Behavioural differences between these two tiers of the same product also 
extend past which search index they seem to ground on.
Citations showed negative selection drift on Plus but
trended positive on Business
(Section~\ref{sec:selection-drift}). Reddit appeared frequently as a
resource on Plus but almost never on Business, and news-related
sources were used much more inline on the Business tier
(Section~\ref{sec:invisible-links}). Running the same prompts on the
same model under two deployment configurations gave us a natural
contrast that strengthened the analysis. 

This thesis opened with the observation that commercial queries
trigger web search at 3$\times$ the overall baseline rate
(Figure~\ref{fig:trigger-rate}), making product recommendations the
arena where retrieval is most active and where source selection is
most consequential. Our data supports the view that retrieval
dependency is the foundation of AI visibility: without SERP ranking,
most citations do not occur.

As this dependence on search indices deepens, so does the contest
over access to them. OpenAI has been accused of scraping Google's
index \parencite{theinformation2025openaiserpapi}, Google has
responded with litigation against third-party scraping services
\parencite{google2025serpapi}, Reddit is suing those same services
over data access \parencite{reddit2025perplexity}, and Microsoft has retired the Bing Search API
entirely (Section~\ref{sec:structural-pivot}). The programmatic access
to search indices is caught in the middle of this access war. The empirical data documented in
this thesis captures a snapshot of a system in rapid flux;
replicating or extending this work will require either new access
agreements with search providers or increasingly creative
instrumentation.

The trajectory points to more retrieval, not less. Semantic fidelity
is already high when using RAG for commercial queries, and commercial
domains continuously produce fresh content for models to ground on.
Retrieval is also far cheaper than inference
(Section~\ref{sec:rag-economics}), giving model providers an economic
reason to lean on search rather than scale up parametric knowledge.
If today's search trigger rates represent the floor rather than the
ceiling, SEO and GEO will only align further. In our data, GEO is
best understood as an extension of SEO rather than a discipline of
its own.
